%% ============================================================
%% 分片 08：原文第 4 节「Compactification And The Geometry Of Hitchin's Moduli Space」（原文 glshort.tex 起于第 2325 行）
%% 由 tools/harvmac2tex.py 从 arXiv 源码（hep-th/0604151v3）自动生成。
%% 只译英文散文；公式、编号、\label、\cite 一律不改。
%% ============================================================

\section{Compactification And The Geometry Of Hitchin's Moduli Space}
\def\MH{{\EUM}_H}

\label{sec:compacthitchin}

We now consider compactification to two dimensions.  We take our
four-manifold to be $M=\Sigma\times C$, a product of two
two-manifolds.  We take $C$ to be a compact Riemann surface on
which we will study the geometric Langlands program.  Generally,
we assume that the genus of $C$ is greater than one, so that for
simple gauge group $G$, the generic flat connection on $C$ is
irreducible. $\Sigma$ may be either a complete but noncompact
two-manifold such as $\Bbb{R}^2$, or a second compact Riemann
surface.   In the latter case, we assume that the size of $\Sigma$
is much larger than that of $C$. We aim to reduce the
four-dimensional gauge theory on $M=\Sigma\times C$ to an
alternative description by an effective field theory on $\Sigma$.

To find the effective physics on $\Sigma$, we must find the
configurations on $M$ that minimize or nearly minimize the action
(Euclidean signature) or the energy (Lorentz signature).  Either
way, one arrives at the conclusion obtained in \cite{ref18}, with closely
related results in \cite{ref19}: the four-dimensional twisted ${\EUN}=4$
supersymmetric gauge theory reduces on $\Sigma$ to a sigma-model
of maps $\Phi:\Sigma\to\MH$, where $\MH$ is the moduli space of
solutions on $C$ of Hitchin's equations  with gauge group $G$.

This  follows from our results in section 第~\ref{sec:toplag}~节.  Apart from the
topological term, which does not affect the equations of motion or
the energy of a classical configuration, the bosonic action of the
twisted ${\EUN}=4$ super Yang-Mills theory is a sum of squares. As
we see from \eqref{eq:wingo}\ and \eqref{eq:uwingo}, it is minimized precisely if the
topological equations are obeyed for all $t$, or in other words if
${\cal F}=D^*\phi=0$, along with
$D\sigma=[\phi,\sigma]=[\sigma,\bar\sigma]=0$. One simple type of
solution of these equations is obtained by taking $A$ and $\phi$
to be pullbacks from $C$, with $\sigma=0$. In fact, the equations
${\cal F}=D^*\phi=0$ for a connection $A$ and an adjoint-valued
one-form $\phi$ on the Riemann surface $C$ are one way to present
Hitchin's equations \cite{ref20}.  By taking the real and imaginary
parts of ${\cal F}$, these equations can alternatively be written
% 式 (4.1)
\setcounter{equation}{0}
\begin{equation}\label{eq:zongo}
\begin{aligned} F-\phi\wedge\phi & = 0\\ D\phi=D^*\phi&=0
.\\\end{aligned}
\end{equation}
  We will write $\MH$, or in more detail $\MH(G,C)$, for the
moduli space of solutions of these equations on $C$ up to gauge
transformations.  If the genus $g$ of $C$ is greater than one,
then $\MH(G,C)$ is a complex manifold of dimension $(2g-2){\rm
dim}\,G$.  It has singularities corresponding to reducible
solutions of Hitchin's equations, that is, solutions with
continuous unbroken gauge symmetries.

As we explained at the end of section 第~\ref{sec:vanishing}~节, any solution of
the equations ${\cal F}=D^*\phi=0$ on $M$ is associated with a map
from $\Sigma$ to $\MH$.  Moreover, if the singularities of $\MH$
are avoided, the solution on $M$ is actually a pullback from a
solution on $C$.  In this case, the conditions on $\sigma$ require
$\sigma=0$.
 So a  field of zero
energy or action that avoids singularities of $\MH$ is a constant
map from $\Sigma$ to a smooth point in  $\M_H$.

Accordingly, a field of almost zero energy or action comes from a
slowly varying map from $\Sigma$ to $\M_H$. (A map is slowly
varying if it varies significantly only over a length scale in
$\Sigma$ that is much greater than the size of $C$.)   So, as long
as we do not encounter fields on $C$ with continuous gauge
symmetries, the effective physics at long distances on $\Sigma$ is
given by a supersymmetric sigma-model in which the target space is
$\M_H(G,C)$.

The importance of the assumption that $\Sigma$ is very large
compared to $C$ is that  it ensures  that nonconstant but slowly
varying maps $\Sigma\to \MH(G,C)$ can exist, and moreover that a
pair of fields $(A,\phi)$ that determine a slowly varying map
$\Phi:\Sigma\to \MH(G,C)$ has very small action.\footnote{There is
also a converse, which we explain in section 第~\ref{sec:wilbranes}~节.  Given a
slowly varying map $\Phi:\Sigma\to \MH(G,C)$, one can reconstruct
a pair of fields $(A,\phi)$ on $M$ that determine this map and
have very small action. } If alternatively we assume that $C$ is
much larger than $\Sigma$, a pair $(A,\phi)$ representing a
generic nonconstant map $\Phi:\Sigma\to \MH(G,C)$ has large
action; instead, we could get a good description involving slowly
varying maps $\tilde\Phi:C\to \MH(G,\Sigma)$.

The above analysis assumes that $G$ is semi-simple. If $G=U(1)$ or
$U(N)$,  some modification is required, because every solution of
Hitchin's equations has a $U(1)$ group of symmetries. As a result,
the low energy description involves the product of a sigma-model
and a supersymmetric $U(1)$ gauge theory. Moreover, for any $G$,
the low energy description breaks down at singularities of $\M_H$
corresponding to solutions of Hitchin's equations that have
continuous gauge symmetries (or extra continuous gauge symmetries
for $G=U(1)$ or $U(N)$). At such points, additional degrees of
freedom come into play: a field of zero energy or action can have
$\sigma\not=0$ and need not be a pullback from $C$.  (It has a
more general decomposition $E=\oplus_{i=1}^nE'_i\otimes E''_i$
described at the end of section 第~\ref{sec:vanishing}~节.)  For semi-simple $G$,
the generic solution of Hitchin's equations is irreducible if the
genus of $C$ is at least 2.  In this paper, we will keep away from
the singularities of $\M_H$ (by considering the Langlands
correspondence for irreducible flat $^L\negthinspace G_{\Bbb{C}}$
bundles, for example).  Some simple generalizations will be
considered in \cite{ref26}.

The sigma-model with target $\MH(G,C)$ will prove to be a powerful
tool for understanding the geometric Langlands program.  But its
relation to the underlying gauge theory is important.  For one
thing, key actors in the story are the  Wilson and 't Hooft line
operators, which are most naturally constructed in the underlying
four-dimensional gauge theory. In addition, the underlying gauge
theory is an important tool when the sigma-model breaks down
because of singularities of $\MH$. But many things can be
understood directly in the sigma-model.

The rest of this section is devoted to a review of the classical
geometry of $\MH$, focussing on just those topics that we need for
the present paper.
